\chapter{Turn Sequence and Energy}

Each turn has an allocation segment followed by twelve impulses. Each impulse contains movement, missile launch, and direct-fire steps in that order.

\section{Complete sequence of play}

\begin{enumerate}
  \item \phase{Allocate Energy.} Each commander secretly plans every surviving ship.
  \item \phase{Commit Fleets.} When both commands are ready, reveal allocations and announce each ship's speed.
  \item \phase{Determine Initiative.} Each side rolls one die; the high roll wins. Re-roll ties.
  \item \phase{Resolve Twelve Impulses.}
  \begin{enumerate}[label=\alph*.]
    \item Draw the next card from the current phase deck.
    \item Move every ship whose speed appears, slowest first. Then move all missiles already on the map.
    \item Launch any missiles.
    \item Fire any eligible lance beams and mass drivers.
  \end{enumerate}
  \item \phase{Check Victory.} Resolve destruction and determine whether the battle has ended.
  \item \phase{End the Turn.} Complete funded shield repairs, clear expended allocations, reset unfired status, and begin a new turn.
\end{enumerate}

\section{Allocate energy}

Each undamaged reactor / engine box generates one energy. Secretly assign that energy among:

\begin{itemize}
  \item \term{Speed}: one energy buys one point of speed, from 0 to 12.
  \item \term{Lance beams}: two energy charges one beam weapon.
  \item \term{Mass drivers}: one energy charges one mass-driver weapon.
  \item \term{Shield reinforcement}: one energy absorbs one damage against the chosen shield during this turn.
  \item \term{Shield repair}: two energy schedules one damaged shield box for restoration at the end of the turn.
  \item \term{Weapon repair}: three energy repairs one eligible direct-fire weapon when that optional rule is enabled.
\end{itemize}

Missiles require no energy. A commander is never required to spend all available energy; unspent energy is lost when the allocation is committed.

When commanding multiple ships, complete and save a plan for each surviving ship before committing the fleet. Committing makes every ship's plan final for the turn.

\section{Announce speed and initiative}

After both commands commit, announce every ship's speed. All other allocation details also become public in tabletop play. Each side rolls one die for initiative; high roll wins, with ties re-rolled.

Initiative settles simultaneous choices. When equal-speed ships are due to move, the side with initiative determines the order. If both sides want to use a special maneuver at the same moment, initiative decides which resolves first. It likewise settles ambiguous shield facings and, where necessary, firing declarations.

\section{Impulse decks}

The twelve cards are arranged as three separately shuffled phase decks:

\begin{center}
\begin{tabular}{lccc}\toprule
Impulses & 1--4 & 5--8 & 9--12\\\midrule
Deck & Phase 1 & Phase 2 & Phase 3\\\bottomrule
\end{tabular}
\end{center}

Each card lists the speeds that receive one movement opportunity. A ship moves if its announced speed appears on the card. Across the complete twelve-card sequence, a ship moving at speed $N$ receives exactly $N$ movement opportunities.

\important{SLOWEST FIRST}{Resolve eligible ships from lowest speed to highest speed. When two ships have the same speed, use initiative to determine their order. Every eligible ship receives its own movement choice.}

\section{Worked impulse}

Suppose an impulse card lists speeds 4, 6, 8, 10, and 12. A speed-4 enemy cruiser and two friendly speed-6 frigates are on the map.

\begin{enumerate}
  \item The speed-4 cruiser resolves its movement first.
  \item The two speed-6 frigates then resolve separately in the order determined by initiative. One may move forward while the other side-slips.
  \item All missiles launched in earlier impulses move up to two hexes toward their assigned targets.
  \item Ships may launch new missiles. Those missiles remain in their launch hex during this impulse.
  \item Eligible ships may fire charged lance beams and mass drivers, one weapon at a time.
  \item When both sides finish firing, end the impulse and draw the next card.
\end{enumerate}
